\documentclass[10pt,a4paper]{amsart}
\usepackage[margin=19mm]{geometry}
\usepackage{amsmath,amssymb,mathtools}
\usepackage[scr=rsfs,cal=euler]{mathalfa}
\usepackage{xfrac}
\usepackage[unicode,hidelinks,bookmarks=false]{hyperref}
\newcommand{\ie}{i.e.\ }
\newcommand{\cf}{cf.\ }
\newcommand{\resp}{respectively\ }
\newcommand{\Subsection}{\subsection}
\providecommand{\nobreakdash}{\nobreak\hbox{-}}
\setlength{\emergencystretch}{2em}
\multlinegap0pt
\usepackage{amssymb}
\usepackage{xcolor}
\usepackage{float}
\newtheorem{theorem}{Theorem}
\title{\textbf{On the cumulative residual interval entropy of doubly truncated random variables}}
\author{Stathis Chadjiconstantinidis, Apostolos Bozikas}
\date{Source: arXiv:2603.16037v1 (CC BY 4.0)}
\begin{document}
\maketitle

\noindent\emph{Source and licence.} \textbf{On the cumulative residual interval entropy of doubly truncated random variables}. Version arXiv:2603.16037v1 (CC BY 4.0), distributed by arXiv under the Creative Commons Attribution 4.0 International licence, http://creativecommons.org/licenses/by/4.0/, as recorded in the arXiv licence metadata for this exact version and shown on the abstract page https://arxiv.org/abs/2603.16037v1. Full licence notice: \texttt{licenses/arxiv-2603.16037-CC-BY-4.0.txt}.\par\smallskip

\noindent\textit{Editorial scope.} Counted unit: the article's theorem theorem (source0.tex lines 529--533). The article's complete original proof of it is delivered with it (source0.tex lines 535--572, one \texttt{proof} environment of 38 lines). No mathematics is written, repaired or re-derived: the statement and the proof are the article's own lines, in the article's own notation, and no retained line is substituted or re-flowed.  No further source-local context is needed: every cross-reference of every retained line resolves inside the two delivered spans.  Nothing was omitted from any retained span, so the package carries no omission record.  No retained line contains a citation, so the wrapper carries no bibliography and no bibliography entry.  No retained line contains a figure environment, an image-inclusion command or drawing code, so no image is delivered and the package declares no asset dependency.  Editorial formatting only, in addition: an AMS \texttt{amsart} preamble that restates the article's own declarations and macros used by the retained lines, namely the environments \texttt{theorem} on separate counters, together with the standard packages those lines need.  The retained environments print in this document as Theorem 1 in order of appearance, whereas the article prints the same items with the numbering of its own full document.  The complete native source of the article as distributed in the arXiv e-print for this exact version is bundled unchanged as \texttt{sources/196513/source0.tex}.
\begin{theorem} Let $X$ be a non-negative absolutely continuous random variable. Then for all $(\tau_1,\tau_2) \in D_X$
	\begin{equation}
		\textstyle H(X;\tau_1,\tau_2) = \mathrm{Cov}\left(X, -\ln[\overline{F}_X(X)-\overline{F}_X(\tau_2)] \mid \tau_1 \leq X \leq \tau_2 \right).
	\end{equation}
\end{theorem}

\begin{proof} 
	We have
	\begin{align}
		&\mathrm{Cov}\left(X, -\ln[\overline{F}_X(X)-\overline{F}_X(\tau_2)] \mid \tau_1 \leq X \leq \tau_2 \right) \nonumber \\
		& \textstyle \qquad = -\mathbb{E}[X \ln[\overline{F}_X(X)-\overline{F}_X(\tau_2)] \mid \tau_1 \leq X \leq \tau_2] \nonumber \\
		& \textstyle \qquad \qquad + \mu_X(\tau_1,\tau_2) \mathbb{E}[\ln[\overline{F}_X(X)-\overline{F}_X(\tau_2)] \mid \tau_1 \leq X \leq \tau_2].
	\end{align}
	It holds
	\begin{align}
		& \mathbb{E}[\ln[\overline{F}_X(X)-\overline{F}_X(\tau_2)] \mid \tau_1 \leq X \leq \tau_2] \nonumber \\
		& \textstyle \qquad \qquad \qquad \qquad = \frac{1}{\overline{F}_X(\tau_1)-\overline{F}_X(\tau_2)} \int_{\tau_1}^{\tau_2} \ln[\overline{F}_X(x)-\overline{F}_X(\tau_2)] f_X(x) dx \nonumber \\
		& \textstyle \qquad \qquad \qquad \qquad = \frac{1}{\overline{F}_X(\tau_1)-\overline{F}_X(\tau_2)} \int_{\tau_1}^{\tau_2} \ln[\overline{F}_X(x)-\overline{F}_X(\tau_2)] (\overline{F}_X(x))' dx \nonumber \\
		& \textstyle \qquad \qquad \qquad \qquad = -\frac{1}{\overline{F}_X(\tau_1)-\overline{F}_X(\tau_2)} \int_{\tau_1}^{\tau_2} \ln[\overline{F}_X(x)-\overline{F}_X(\tau_2)] (\overline{F}_X(x)-\overline{F}_X(\tau_2))' dx \nonumber \\
		& \textstyle \qquad \qquad \qquad \qquad = \ln[\overline{F}_X(\tau_1)-\overline{F}_X(\tau_2)] - 1,
	\end{align}
	and similarly,
	\begin{align}
		& \qquad \mathbb{E}[X \ln[\overline{F}_X(X)-\overline{F}_X(\tau_2)] \mid \tau_1 \leq X \leq \tau_2] \nonumber \\
		& \textstyle \qquad \quad = -\frac{1}{\overline{F}_X(\tau_1)-\overline{F}_X(\tau_2)} \int_{\tau_1}^{\tau_2} x \ln[\overline{F}_X(x)-\overline{F}_X(\tau_2)] (\overline{F}_X(x)-\overline{F}_X(\tau_2))' dx \nonumber \\
		& \textstyle \qquad \quad = \tau_1 \ln[\overline{F}_X(\tau_1)-\overline{F}_X(\tau_2)] \nonumber \\
		& \textstyle \qquad \qquad + \frac{1}{\overline{F}_X(\tau_1)-\overline{F}_X(\tau_2)} \int_{\tau_1}^{\tau_2} [\overline{F}_X(x)-\overline{F}_X(\tau_2)] \ln[\overline{F}_X(x)-\overline{F}_X(\tau_2)] dx - \mu_X(\tau_1,\tau_2).
	\end{align}
	Hence, substituting (20) and (21) into the right-hand side of (19) we obtain
	\begin{align}
		& \textstyle \qquad \mathrm{Cov}\left(X, -\ln[\overline{F}_X(X)-\overline{F}_X(\tau_2)] \mid \tau_1 \leq X \leq \tau_2 \right) = -\tau_1 \ln[\overline{F}_X(\tau_1)-\overline{F}_X(\tau_2)] \nonumber \\
		& \textstyle \qquad \quad - \int_{\tau_1}^{\tau_2} \frac{\overline{F}_X(x)-\overline{F}_X(\tau_2)}{\overline{F}_X(\tau_1)-\overline{F}_X(\tau_2)} \ln[\overline{F}_X(x)-\overline{F}_X(\tau_2)] dx \nonumber \\
		& \textstyle \qquad \qquad + \mu_X(\tau_1,\tau_2) \ln[\overline{F}_X(\tau_1)-\overline{F}_X(\tau_2)]. 
	\end{align}
	Now, using (8) and (10) it follows that
	$$ \textstyle \mu_X(\tau_1, \tau_2) = \tau_1 + \int_{\tau_1}^{\tau_2} \frac{\overline{F}_X(x) - \overline{F}_X(\tau_2)}{\overline{F}_X(\tau_1) - \overline{F}_X(\tau_2)} \, dx, $$
	and substituting this into the right-hand side of (22), we get
	\begin{align*}
		& \textstyle \qquad \mathrm{Cov}\left(X, -\ln\left[\overline{F}_X(X) - \overline{F}_X(\tau_2)\right] \,\middle|\, \tau_1 \leq X \leq \tau_2 \right) \\
		& \textstyle \qquad \quad = -\int_{\tau_1}^{\tau_2} \frac{\overline{F}_X(x) - \overline{F}_X(\tau_2)}{\overline{F}_X(\tau_1) - \overline{F}_X(\tau_2)} \ln\left[\overline{F}_X(x) - \overline{F}_X(\tau_2)\right] \, dx \\
		& \textstyle \qquad \qquad + \ln\left[\overline{F}_X(\tau_1) - \overline{F}_X(\tau_2)\right] \int_{\tau_1}^{\tau_2} \frac{\overline{F}_X(x) - \overline{F}_X(\tau_2)}{\overline{F}_X(\tau_1) - \overline{F}_X(\tau_2)} \, dx \\
		& \textstyle \qquad \quad = H(X; \tau_1, \tau_2).
	\end{align*}
\end{proof}

\end{document}
